\scanpage{175}\mspage{30}%
phrase an extremely simple topological intuition, which I am going to make explicit, assuming to fix ideas that we are over the base field $\mathbb C$ and that we take advantage of this to work with
% Compact isolated point at whole-image (171,423), left of L5; origin and role uncertain.
\raisebox{0pt}[0pt][0pt]{\parbox{0pt}{\figure{\begin{tikzpicture}[overlay,x=1pt,y=1pt]\fill (-7,0) ellipse[x radius=.35pt,y radius=.5pt];\end{tikzpicture}}}}%
the usual locally compact spaces. The topological idea behind the equisingularity hypothesis is that there exists a \uline{tubular neighborhood $T$} of $X_0$ in $X$ retracting onto $X_0$, and such that $(X_0,T)$ over $X_0$ is a \unsure{locally|locally} trivial bundle, hence $T-X_0$ over $X_0$ is locally trivial. In fact, \deleted{\unsure{we shall be able to con|we shall be able to cons}} \inserted{if $\partial T$ is the} ``boundary'' of $T$, which will also be a locally trivial bundle over $X_0$, $T$ over $X$ is the induced cone bundle ($\simeq(\partial T\times I)\amalg_{\partial T}X_0$, where $I=[0,1]$, $\partial T%
% Compact isolated point at whole-image (702,1183), above the first relation shaft; not a mathematical dot.
\raisebox{1.3ex}[0pt][0pt]{\parbox{0pt}{\figure{\begin{tikzpicture}[overlay,x=1pt,y=1pt]\fill (4,0) -- (4.7,.6) -- (4.8,-.5) -- cycle;\end{tikzpicture}}}}%
\mathrel{\to}\partial T\times I$ defined by $x\mathrel{\unsure{\mapsto|\vdash}}(x,1)$, and $\partial T\mathrel{\unsure{-|\to}}X_0$ is the projection) hence $\mathring T=T-X_0\simeq\partial T\times[0,1[$ is $X_0$-homotopic to $\partial T$. If $X_0$ and $X_1$ are nonsingular, the $\unsure{=|\overline m}$ holds for $\mathring T$ and $\partial T$, which are topologically smooth bundles over $X_0$. Moreover, $\deleted{\unsure{X_1|X}}$ putting $\widetilde X_1=X_1-\mathring T$, \deleted{$\unsure{\widetilde X_1|X_1}$ dev.} the inclusion $\widetilde X_1%
% Short isolated dark tick at whole-image (259,1797), below this source object; origin and role uncertain.
\raisebox{-2.3ex}[0pt][0pt]{\parbox{0pt}{\figure{\begin{tikzpicture}[overlay,x=1pt,y=1pt]\draw[line width=.55pt] (0,.6) -- (.2,-.6);\end{tikzpicture}}}}%
\hookrightarrow X_1$ is a homotopy equivalence, and $X$ is reconstructed, up to homeomorphism, from \unsure{the|the} diagram of spaces
\diagram{\begin{tikzcd}[ampersand replacement=\&,column sep=small,row sep=large]
(\mathring T=T-X_0\underline{\approx}) \& \partial T \arrow[r,"\unsure{J|j}","\text{inclusion}"'] \arrow[d,"p","\text{fibration}"'] \& \widetilde X_1\,(\underline{\approx}X_1)\\
\& X_0 \&
\end{tikzcd}}
as the amalgamated sum. In terms of this diagram of spaces; the functor $\varphi$ more
% archive sheet number: 175
